\documentclass[11pt]{article}
\usepackage[margin=1in]{geometry}
\usepackage{enumitem}
\usepackage{xcolor}
\usepackage{tcolorbox}
\usepackage{booktabs}
\usepackage{tabularx}
\usepackage{array}
\usepackage{hyperref}

\definecolor{required}{HTML}{c0392b}
\definecolor{easy}{HTML}{27ae60}
\definecolor{medium}{HTML}{e67e22}
\definecolor{quoteblue}{HTML}{2c3e50}
\definecolor{boxgray}{HTML}{f5f5f5}
\definecolor{chatbg}{HTML}{eef2f7}
\definecolor{chatframe}{HTML}{5b7ea1}

\newcommand{\tagR}{\textcolor{required}{\textbf{[REQUIRED]}}}
\newcommand{\tagE}{\textcolor{easy}{\textbf{[EASY]}}}
\newcommand{\tagM}{\textcolor{medium}{\textbf{[MEDIUM]}}}
\newcolumntype{Y}{>{\raggedright\arraybackslash}X}
\newcolumntype{C}[1]{>{\centering\arraybackslash}p{#1}}

\newtcolorbox{reviewquote}{
  colback=boxgray, colframe=quoteblue,
  boxrule=0.5pt, leftrule=3pt,
  arc=0pt, outer arc=0pt,
  top=6pt, bottom=6pt, left=10pt, right=10pt
}

\newtcolorbox{chatquote}{
  colback=chatbg, colframe=chatframe,
  boxrule=0.5pt, leftrule=3pt,
  arc=0pt, outer arc=0pt,
  top=6pt, bottom=6pt, left=10pt, right=10pt
}

\title{\textbf{TOMPECS Review Summary \& Revision Plan}\\[4pt]
\normalsize Manuscript ID: TOMPECS-2024-0030.R1\\
``Impact of Network Topology on Byzantine Resilience in DFL''}
\author{}
\date{}

\begin{document}
\maketitle
\vspace{-2em}

% ============================================================
\part*{I. Raw Reviewer Feedback}
% ============================================================

\section{Associate Editor}

\begin{reviewquote}
The authors have addressed some concerns from the previous round of review, but several important concerns remain. In particular, the manuscript as written does not make clear which results are ``surprising'' and provide insights beyond what is already known or expected from the existing literature on theoretical or experimental investigations of Byzantine resilience in multi-agent optimization. Pointing these out, or conducting more experiments that reveal such unexpected insights, would greatly strengthen the paper. For example, federated learning is a special case of distributed multi-agent optimization---why would we expect it to behave differently under Byzantine attacks? I agree with Reviewer 1 that the author response, which points out the high dimensionality of machine learning, does not sufficiently address this question.

\medskip

The reviewers also point out that several limitations of the paper, e.g., considering only the simple MNIST dataset and two types of network topologies, have not been addressed. I also concur with the reviewers that showing the range of experimental results would be beneficial; even showing evidence that they are tightly clustered around the mean, as claimed in the response, would be useful information to show the reader.
\end{reviewquote}

\section{Reviewer 1}

\begin{reviewquote}
Most of my concerns remain. Vector-input optimization and high-dimensional ML are essentially the same unless the scaling of the dimensions is clearly characterized which is not considered here. The experiments added, i.e., sign-flip attack and trimmed means, those are very basic. Hence, the experiments are still limited. The reviewer did agree that MNIST is used as one of the baselines. Yet, experiments based on MNIST only are not sufficient. At least CIFAR10 should be considered.
\end{reviewquote}

\section{Reviewer 2}

\begin{reviewquote}
Thank you for your response and efforts to address the concerns raised in the previous review comments. Here are some follow-up comments I have on the latest version:

\medskip
\textbf{1.} ``We find that multiple runs of the same experimental setup are typically identical, \ldots''

\medskip
I still think, given that this paper is an experiment-oriented study, it would be better to include as many experimental results as possible. Since the authors have already conducted multiple runs under the same setups, such results can be aggregated into a single plot, for example, showing the average accuracy and standard deviation. It would be better, in my opinion, to present those results instead of omitting them, to make the paper more convincing.

\medskip
As a side note, I assume that multiple runs of the same setup you've mentioned would have used different randomness seeds to generate, e.g., the network, the partition of the dataset, the initial estimate, the computation of stochastic gradients, etc, so their results could be similar but definitely not identical.

\medskip
\textbf{2.} ``It is true that failure of convergence is expected when over half of the nodes in the network are Byzantine \ldots''

\medskip
The reason I mentioned this issue, that it is already expected to not converge when Byzantine proportion $>$50\%, is that these results do not add to the paper's contribution. From my perspective, it would be more interesting to see the differences in the actual performance of the algorithms when facing different proportions of Byzantine agents, e.g., 10\%, 20\%, 30\%, 40\%, \ldots, instead of confirming that the algorithms will indeed fail when $>$50\%. I think the paper would contribute more if the authors could add some experiments showing this kind of results.

\medskip
It is worth noting that even some theory-oriented papers on this topic have tested different failure proportions in their experiment sections before, e.g., Figures 5 and 6 in [A] and Figure 2 in [B]. It is expected that an experiment-oriented paper will do a more comprehensive study.

\medskip
\textbf{Typos:}
\begin{itemize}[nosep]
  \item Line 251: ``\ldots determine the degree of each node to be \textit{4}.''
  \item Line 435: Fig.~3(a) ``\textit{AAccuracy} of Krum \ldots''
\end{itemize}

\medskip
\textbf{Ref.}\\
{[A]} Blanchard, Peva, et al. ``Machine learning with adversaries: Byzantine tolerant gradient descent.'' NeurIPS 2017.\\
{[B]} Liu, Shuo, Nirupam Gupta, and Nitin H. Vaidya. ``Impact of redundancy on resilience in distributed optimization and learning.'' ICDCN 2023.
\end{reviewquote}

% ============================================================
\newpage
\part*{II. Claude Chat Analysis}
% ============================================================

\section{Initial Assessment}

\begin{chatquote}
Okay this is actually great news --- these are extremely actionable comments. This is not a ``fundamental contribution is weak'' rejection. The reviewers know exactly what they want and it's all experiments + framing.
\end{chatquote}

\section{The Three Real Asks}

\begin{chatquote}
\textbf{1. Add CIFAR-10} (Reviewer 1, AE). Non-negotiable, this is the clearest single request. Rerun your existing experiment pipeline with CIFAR-10 instead of MNIST. Same topologies, same aggregation functions. Claude Code + the existing GitHub repo makes this maybe a weekend of compute time.

\medskip
\textbf{2. Add error bars / std dev across runs} (Reviewer 2, point 1). They literally told you they'd accept just showing the runs are tightly clustered. Run each setup 3--5 times with different seeds, plot mean $\pm$ std. This is a few lines of matplotlib.

\medskip
\textbf{3. Finer Byzantine proportions} (Reviewer 2, point 2). Instead of 0/30/60\%, test 10/20/30/40/50\%. They explicitly want to see the degradation curve, not just the cliff edge. Again, just rerunning with different parameters.
\end{chatquote}

\section{The Framing Ask}

\begin{chatquote}
The framing ask (AE) is the only intellectually harder one --- they want you to articulate why FL topology matters differently than generic distributed optimization. You'd need to write 1--2 paragraphs in the intro/discussion making this case. The high-dimensionality argument you used before didn't land --- you'd need a different angle, probably around the P2P dynamic topology argument you already have in the paper.
\end{chatquote}

\section{Typos}

\begin{chatquote}
The two typos --- Line 251 and the ``AAccuracy'' --- are embarrassingly easy fixes.
\end{chatquote}

\section{Bottom Line}

\begin{chatquote}
Honestl, this revision is like 80\% compute jobs and 20\% writing. With Claude Code handling the experiment reruns, this is genuinely a few weekends of work. The journal is still open to it --- the deadline expired but that's recoverable by emailing Golubchik directly explaining the situation.
\end{chatquote}

\section{Strategic Concern (18-Month Gap)}

\begin{chatquote}
\textbf{What's likely changed since April 2024:}

The decentralized FL + Byzantine robustness space has gotten more crowded. There's been a lot of work on robust aggregation in non-IID settings, and some papers have probably touched adjacent topology questions. You'd need to do a literature sweep to check if anyone published something that partially scoops your contribution.

\medskip
\textbf{The journal situation:}

The submission is technically dead --- deadline expired January 2025. But TOMPECS is not a fast-moving venue and Golubchik is the EIC. A polite email explaining ``the lead author graduated and I'm reviving this as first author'' is genuinely recoverable. Editors deal with this. The worst she says is ``resubmit as new.''

\medskip
\textbf{My honest take though:}

Given the time gap, I'd actually lean toward a different strategy than just reviving this paper as-is for TOMPECS. Consider:
\begin{enumerate}[nosep]
  \item Do the CIFAR-10 experiments and finer Byzantine proportions anyway
  \item Update the related work to cover 18 months of new papers
  \item Reframe around whatever the most current angle is
  \item Submit to a workshop at NeurIPS or ICML 2025 instead
\end{enumerate}

\medskip
A workshop paper has faster turnaround, lower bar, and frankly better visibility in the ML community than TOMPECS --- which is more of a systems/networking venue that Stanford AI faculty probably don't read.

\medskip
The experiments you need to add are the same either way. The destination is the question.
\end{chatquote}

\end{document}